\chapter*{Preface: Why Set the Textbooks Alight?}
\addcontentsline{toc}{chapter}{Preface: Why Set the Textbooks Alight?}
\markboth{Preface}{}

First, let's be clear: ``igniting'' in the title does not mean actually lighting a fire. Although burning paper leaves ash and releases smoke, that is a chemical reaction for another chapter. What we want to ignite is your curiosity about the world. What we want to burn away is a particular way of learning chemistry and biology: memorization.

Memorize element symbols, memorize equations, memorize ``mitochondria are the cell's powerhouses,'' memorize ``dominant genes control dominant traits.'' Finish one exam, then start memorizing for the next. If you have ever stared at the periodic table and wondered what any of it has to do with you, this book is for you.

The bread you eat each day, the scab over a cut on your hand, the apple that slowly browns after you slice it, and the gas that makes soda fizz are all performances of the same laws at different scales. Those laws can be arranged in a long chain of cause and effect:

\begin{quote}
Elementary particles \(\rightarrow\) atoms \(\rightarrow\) periodic table \(\rightarrow\) chemical bonds \(\rightarrow\) molecular structure \(\rightarrow\) chemical reactions \(\rightarrow\) organic molecules \(\rightarrow\) biological macromolecules \(\rightarrow\) cells \(\rightarrow\) metabolism \(\rightarrow\) DNA \(\rightarrow\) gene expression \(\rightarrow\) inheritance \(\rightarrow\) evolution.
\end{quote}

School cuts this chain in two, calls the first half chemistry and the second biology, and then cuts it again into chapters of facts. This book does the opposite: we follow the chain from beginning to end. You will discover that life is no mysterious force beyond chemistry. It consists of ordinary molecules organized to an extraordinary degree, able to preserve information from the past, maintain their own order, and create a future that will continue to change.

After the journey, you should be able to do things you could not do before: explain the periodic table's shape through electron structure instead of memorizing groups and periods; use energy and probability to judge whether a reaction can happen and how fast; understand what the molecular skeletons of aspirin and caffeine tell you; distinguish DNA, genes, chromosomes, and genomes; and judge claims about cancer-curing supplements, genetic fortune-telling, and genetic-modification horror stories for yourself.

We cannot promise every page will be easy. Some reasoning requires you to stop, pick up a pen, draw, and calculate. But we can promise three things. First, no skipped steps: each conclusion grows from something you already understand. Second, no pretending: when a question remains unanswered, such as precisely how life originated, we will frankly say we do not know. Third, no pointless memorization: before each symbol appears, you will encounter the real phenomenon behind it.

Once curiosity catches fire, the dry sentences in a textbook become living stories. Now, let the spark catch.

\hfill The Author
